Organisations increasingly derive emotion estimates from live video. A deployed affect recognition component must stay within a compute budget, cope with noise, compression and poor lighting without labels, and remain accountable. Following design science research, we developed a multimodal affect recognition artifact over six build-and-evaluate cycles, on 1440 acted clips from 24 actors, and confirmed the findings on 2178 clips from 91 further actors and on 2329 spontaneous clips from 1254 online videos, with unseen people or videos. Late fusion of prosodic, wav2vec~2.0, MobileNet and CLIP features reached 82.6\% accuracy over eight acted emotions and 55.8\% over four spontaneous ones, 5.5--10.5 points above the best single modality. Running the speech model first in a confidence-gated cascade cut CPU cost by 38--48\%. Under audio noise, clean-trained models became more confident as their accuracy fell, silently disabling the cascade and confidence-based drift detection. Label-free per-channel monitors of signal quality detected every degradation in our acted test streams and raised fewer false alarms on speaker changes than content-based monitors. In online video, background noise and low bit rates are part of the clean reference, so the monitors raised almost no false alarms but missed such degradations. Dropping a flagged channel restored accuracy under 10\,dB noise from 37\% to 61\% on acted data. Training on degraded data helped more, but it turned gating on clean-data alarms harmful. Models trained on acted data were confidently wrong on online video. We derive six design principles centred on measured cost, robust training, quality monitoring and envelope-aware gating.
